% Generated draft: requires mathematical and source review.
\documentclass[11pt]{article}
\usepackage[margin=25mm]{geometry}
\usepackage{fontspec}
\setmainfont{Latin Modern Roman}
\usepackage{amsmath,amssymb,mathtools,unicode-math}
\setmathfont{Latin Modern Math}
\usepackage{enumitem,needspace,xurl,hyperref}
\setlength{\parindent}{0pt}
\setlength{\parskip}{5pt}
\setlength{\emergencystretch}{3em}
\newcommand{\N}{\mathbb{N}}
\newcommand{\Z}{\mathbb{Z}}
\newcommand{\Q}{\mathbb{Q}}
\newcommand{\R}{\mathbb{R}}
\newcommand{\C}{\mathbb{C}}
\newcommand{\F}{\mathbb{F}}
\newcommand{\A}{\mathbb{A}}
\newcommand{\PP}{\mathbb{P}}
\newcommand{\E}{\mathbb{E}}
\newcommand{\eps}{\varepsilon}
\newcommand{\dd}{\,\mathrm{d}}
\begin{document}
\section*{UnsolvedMath ID: 20001642}
\textbf{Source problem code:} catalogue entry
\subsection{Definitions and mathematical statement}
Let \(n\geq2\) and let \(CV_n\) be Culler--Vogtmann Outer space of marked volume-one metric core graphs of rank \(n\). Use the directed Lipschitz metric
\[
 d_L(X,Y)=\log\inf\{\operatorname{Lip}(f):f:X\to Y
 \text{ respects the marking up to homotopy}\}.
\]
It is asymmetric. For distance comparisons one may specify the symmetrization \(d_{\mathrm{sym}}(X,Y)=\max\{d_L(X,Y),d_L(Y,X)\}\); the original source does not choose a symmetrization. An iwip is a fully irreducible outer automorphism: no positive power preserves a proper nontrivial free-factor conjugacy class. Put
\[
 \tau(\phi)=\inf_{Y\in CV_n}d_L(Y,\phi Y),\qquad
 \operatorname{Min}(\phi)=\{X:d_L(X,\phi X)=\tau(\phi)\}.
\]
The source is a five-part program comparing Lipschitz Outer-space geodesics with Teichm\"uller and Weil--Petersson behavior:
\begin{enumerate}
\item Characterize geodesics or lines with strongly contracting projection functions, with specified uniform quantitative constants.
\item Find a suitable thick-geodesic or thick-part notion with hyperbolic behavior. A mere lower bound on graph injectivity radius is explicitly said to be insufficient; ask also about thick-to-thin geodesics and quantitative thickness of every iwip axis.
\item Relate geodesic behavior to an appropriate boundary theory.
\item Determine the geometry and topology of \(\operatorname{Min}(\phi)\): is it CAT(0), contained in a CAT(0) space or an axis bundle, and quasi-isometric or homotopy equivalent to a line? How does it change under powers?
\item Quantify the proximity of \(\operatorname{Min}(\phi)\) and \(\operatorname{Min}(\phi^{-1})\).
\end{enumerate}
The intended directed projection, contraction constants, boundary and metric on subsets are not fully defined. These are research questions, not one single conjecture.

The workshop discussion adds material qualifications. Juan favors thick geodesics over a thick subset because of parallel-ray and stratum phenomena, and asks how contraction varies between iwips. Yael Algom-Kfir describes quantitative axis contraction. Yair Minsky suggests the boundary question and asks about powers. Kasra Rafi proposes a bounded-local-projection heuristic; Juan suggests defining the relevant projections from contraction instead. Handel discusses axes, train tracks for powers, endpoint-defined axis bundles and variable transverse diameter; Mladen Bestvina warns about illegal turns. Karen Vogtmann asks for the set's geometry. Thierry Coulbois questions the specificity of the inverse-minset question, and Martin Lustig tentatively raises a Masur-criterion analogue for non-uniquely-ergodic points. These are attributed discussion remarks, not all established statements.

The aid reports modern results characterizing nondegenerate strongly contracting directed geodesics by parametrized quasigeodesic progress in the free-factor complex. It also reports that an iwip minset is thick, connected, locally finite and cocompact under the cyclic action, is a quasi-line, and has finite symmetric Hausdorff distance from the inverse minset. Here thickness would mean a positive lower bound, depending on \(\phi\), on nontrivial conjugacy-class lengths. The complete nondegeneracy conditions, metric conventions, uniform constants and literature references for these claims are absent; no bound independent of \(\phi\) or CAT(0) conclusion is asserted.

The aid's concrete metric statement uses the iwip translation identities \(\tau(\phi)=\log\lambda_\phi>0\) and \(\tau(\phi^k)=k\tau(\phi)\). It reports that every \(X\in\operatorname{Min}(\phi)\) lies on a \(\phi\)-periodic directed axis obtained by concatenating translates of a minimizing segment. With directed arc-length parameter, such an axis satisfies
\[
 d_L(\gamma(s),\gamma(t))=t-s\ (s<t),\qquad
 \gamma(t+\tau(\phi))=\phi\gamma(t).
\]
For every integer \(k\geq1\), the reported inclusions and exact test are
\[
 \operatorname{Min}(\phi)\subseteq\operatorname{Min}(\phi^k),\qquad
 \Delta_k(X)=k\,d_L(X,\phi X)-d_L(X,\phi^kX)\geq0,
\]
\[
 X\in\operatorname{Min}(\phi^k):\qquad
 \Delta_k(X)=0\quad\Longleftrightarrow\quad X\in\operatorname{Min}(\phi).
\]
For nongeometric iwips the aid invokes a Handel--Mosher axis-bundle decomposition by the union of power minsets. Its exact definition, and whether closure or other qualifications occur in the published formulation, need checking. Inclusions of power minsets are naturally ordered by divisibility, \(\operatorname{Min}(\phi^k)\subseteq\operatorname{Min}(\phi^{k\ell})\), rather than an unqualified claim of nesting for every successive integer. This shortcut-defect criterion is a partial metric description, not a resolution of every workshop subquestion.
\subsection{Short English statement}
The workshop asks five linked questions about contraction, thickness, boundaries and iwip minsets in asymmetric Lipschitz Outer space. The aid reports modern minset and contraction results, and gives a precise power-minset shortcut defect detecting base minpoints. Metric conventions, quantitative hypotheses and the exact axis-bundle statement remain review issues.
\subsection{Sources}
\begin{itemize}
\item Local UnsolvedMath registry, permanent ID \(20001642\), code catalogue entry; canonical \texttt{aim-geometric-group-theory-notes.json}, \texttt{notes[150]}.
\item AIM workshop notes ``The geometry of the outer automorphism group of a free group,'' 25--29 October 2010, edited by Johanna Mangahas, Problem 2.1, \url{https://aimath.org/WWN/outerauto/outerauto.pdf}. The supplied verification gives this identification; no PDF inspection is claimed here.
\item Participants named in the source discussion: Juan (surname not supplied), Yael Algom-Kfir, Yair Minsky, Kasra Rafi, Handel (first name not supplied), Mladen Bestvina, Karen Vogtmann, Thierry Coulbois and Martin Lustig. Their questions, heuristics and uncertainty are preserved as discussion attribution.
\item Handel--Mosher axis bundles and modern contraction/minset theorems are invoked by the aid without titles, dates, theorem numbers or full authoritative citations.
\item Research attempt 1: ``partial\_result,'' ``partially\_solved,'' rigor label ``proved,'' novelty confidence ``low.'' The aid is machine generated, not peer reviewed; independently verify before citation.
\end{itemize}
\end{document}
